\section*{Chapitre \ref{ch:g10:lines} --- Équations de droites et systèmes linéaires}

\begin{solution}{exo:g10:lines:1}
$y = 2x - 3$~: coefficient $2$, \omterm{def:g10:reffunc:affine}{ordonnée à l'origine} $-3$.
$y = -\frac13 x + 1$~: coefficient $-\frac13$, \omterm{def:g10:reffunc:affine}{ordonnée à l'origine} $1$.
$y = 4$~: coefficient $0$, \omterm{def:g10:reffunc:affine}{ordonnée à l'origine} $4$ (droite
horizontale).
$x = -2$~: droite verticale, pas de \omterm{def:g10:reffunc:affine}{coefficient directeur} ni d'\omterm{thm:g10:lines:reduced}{équation
réduite}.
\end{solution}

\begin{solution}{exo:g10:lines:2}
$3 \times 2 + 1 = 7$~: oui, $A$ est sur la droite.
$3 \times (-1) + 1 = -2 \neq -1$~: $B$ n'y est pas.
$3 \times 4 + 1 = 13$~: oui, $C$ est sur la droite.
\end{solution}

\begin{solution}{exo:g10:lines:3}
(a)~$m = \frac{8 - 2}{3 - 0} = 2$ et $p = y_A = 2$ (le point $A$ est
sur l'axe des $y$)~: $y = 2x + 2$.

(b)~$m = \frac{-4 - 5}{2 - (-1)} = \frac{-9}{3} = -3$~; puis
$5 = -3 \times (-1) + p$ donne $p = 2$~: $y = -3x + 2$.

(c)~$x_A = x_B = 4$~: droite verticale $x = 4$.
\end{solution}

\begin{solution}{exo:g10:lines:4}
$y = \frac{6x+2}{2} = 3x + 1$. Les droites de coefficient $3$ sont
$y = 3x - 1$, $y = 3x + 5$ et $y = 3x + 1$~: ces trois-là sont parallèles
entre elles~; aucune paire n'est confondue (\omterm{def:g10:reffunc:affine}{ordonnées à l'origine}
différentes), et aucune n'est parallèle à $y = -3x + 2$.
\end{solution}

\begin{solution}{exo:g10:lines:5}
Premier système~: substituer $y = 2x - 1$ dans $3x + y = 9$~:
$3x + 2x - 1 = 9$, donc $5x = 10$, $x = 2$, puis $y = 3$. Solution~:
$(2, 3)$.

Second système~: de $x - y = 4$, $x = y + 4$. Puis
$2(y + 4) + 3y = 3$, donc $5y = -5$, $y = -1$, puis $x = 3$. Solution~:
$(3, -1)$.
\end{solution}

\begin{solution}{exo:g10:lines:6}
Premier système~: additionner les \omterm{def:g10:algebra:equation}{équations} élimine $y$~: $2x = 14$,
$x = 7$~; puis $7 + y = 10$ donne $y = 3$. Solution~: $(7, 3)$.

Second système~: multiplier la première \omterm{def:g10:algebra:equation}{équation} par $5$ et la seconde
par $-2$~:
\[
\begin{cases}
15x + 10y = 60 \\
-4x - 10y = -16
\end{cases}
\]
Addition~: $11x = 44$, donc $x = 4$~; puis $3 \times 4 + 2y = 12$
donne $y = 0$. Solution~: $(4, 0)$.
\end{solution}

\begin{solution}{exo:g10:lines:7}
Calculer $ab' - a'b$ dans chaque cas.

Premier~: $2 \times 2 - (-4) \times (-1) = 4 - 4 = 0$, et la seconde
\omterm{def:g10:algebra:equation}{équation} est $-2$ fois la première~: deux fois la même droite, une
infinité de solutions.

Second~: $ab' - a'b = 0$ encore, mais les seconds membres ne sont pas
dans le rapport $-2$ ($1 \neq -6$)~: deux droites parallèles distinctes,
aucune solution.

Troisième~: $2 \times 1 - 1 \times (-1) = 3 \neq 0$~: exactement une
solution.
\end{solution}

\begin{solution}{exo:g10:lines:8}
Soit $x$ le nombre de billets enfants et $y$ le nombre de billets
adultes~:
\[
\begin{cases}
x + y = 200 \\
5x + 9y = 1352 .
\end{cases}
\]
De la première \omterm{def:g10:algebra:equation}{équation} $x = 200 - y$~; en substituant,
$5(200 - y) + 9y = 1352$, donc $1000 + 4y = 1352$, $y = 88$, puis
$x = 112$. Donc $112$ billets enfants et $88$ billets adultes.
Vérification~: $5 \times 112 + 9 \times 88 = 560 + 792 = 1352$.
\end{solution}

\begin{solution}{exo:g10:lines:9}
\emph{1.} Parallèle signifie même coefficient $2$~: $y = 2x + p$ avec
$4 = 2 \times 1 + p$, donc $p = 2$~: $d'\colon y = 2x + 2$.

\emph{2.} Sur l'axe des $x$, $y = 0$~: $2x + 2 = 0$ donne $x = -1$. Le
point d'\omterm{def:g10:numbers:interunion}{intersection} est $(-1, 0)$.
\end{solution}

\begin{solution}{exo:g10:lines:10}
\emph{1.} \omterm{prop:g10:coordgeom:midpoint}{Milieu} de $[BC]$~: $A' = \left(\frac{6+2}{2},
\frac{2+4}{2}\right) = (4, 3)$. La médiane issue de $A(0,0)$ a pour
coefficient $\frac{3 - 0}{4 - 0} = \frac34$~: \omterm{def:g10:algebra:equation}{équation}
$y = \frac34 x$.

\omterm{prop:g10:coordgeom:midpoint}{Milieu} de $[AC]$~: $B' = (1, 2)$. La médiane issue de $B(6, 2)$ a pour
coefficient $\frac{2 - 2}{1 - 6} = 0$~: c'est la droite horizontale
$y = 2$.

\emph{2.} \omterm{def:g10:numbers:interunion}{Intersection}~: $\frac34 x = 2$ donne $x = \frac83$, donc
$G\left(\frac83, 2\right)$. La troisième médiane joint $C(2,4)$ au
\omterm{prop:g10:coordgeom:midpoint}{milieu} de $[AB]$, $C' = (3, 1)$~; son coefficient est
$\frac{1 - 4}{3 - 2} = -3$, \omterm{def:g10:algebra:equation}{équation} $y = -3x + 10$. En $x = \frac83$~:
$y = -8 + 10 = 2$~: $G$ est dessus. Et en effet
$G = \left(\frac{0 + 6 + 2}{3}, \frac{0 + 2 + 4}{3}\right)$~: le
centre de gravité moyenne les \omterm{def:g10:coordgeom:system}{coordonnées} des sommets.
\end{solution}

\begin{solution}{pb:g10:lines:1}
\textbf{1.} \omterm{def:g10:reffunc:affine}{Coefficient directeur} $\frac{8 - 2}{3 - 1} = 3$~; par
$(1, 2)$~: $y = 3x - 1$.

\textbf{2.} $y = 3x - 1$ et $y = 3x + 4$~: même \omterm{def:g10:reffunc:affine}{coefficient
directeur}, droites parallèles (et distinctes). En coupant
$y = 3x - 1$ et $y = -x + 7$~: $3x - 1 = -x + 7$ donne $x = 2$,
$y = 5$~: point $(2, 5)$.

\textbf{3.} Les deux points ont la même \omterm{def:g10:coordgeom:system}{abscisse} $x = 2$~: la droite
est verticale, d'\omterm{def:g10:algebra:equation}{équation} $x = 2$. Une forme réduite $y = mx + p$
répond «~un seul $y$ par $x$~», ce que refuse une droite verticale~:
son \omterm{def:g10:reffunc:affine}{coefficient directeur} serait infini.

\textbf{4.} $3 \times 4 - 1 = 11$~: oui, $(4, 11)$ est sur la
droite. $3 \times (-1) - 1 = -4 \neq -5$~: non.

\textbf{5.} $y = -2x + 3$~; la droite coupe les axes en $(0, 3)$ et
$\left(\frac32, 0\right)$. Aire du triangle~:
$\frac12 \times \frac32 \times 3 = \frac94 = 2.25$.

\textbf{6.} Par substitution~: $3x + 2x - 3 = 7$, donc $x = 2$ et
$y = 1$~: solution $(2, 1)$.

\textbf{7.} En additionnant~: $7x = 7$, $x = 1$~; puis $3y = 6$,
$y = 2$~: solution $(1, 2)$.

\textbf{8.} Premier système~: la seconde \omterm{def:g10:algebra:equation}{équation} est le double de la
première --- une droite comptée deux fois, donc une infinité de
solutions (tous les points de $2x - y = 3$). Second système~: mêmes
premiers membres proportionnels, mais pas les seconds
($2 \times 3 = 6 \neq 1$)~: deux droites parallèles --- aucune
solution. Trichotomie~: \omterm{def:g10:reffunc:affine}{coefficients directeurs} distincts $\to$ une
\omterm{def:g10:numbers:interunion}{intersection}~; coefficients égaux et \omterm{def:g10:reffunc:affine}{ordonnées à l'origine}
différentes $\to$ parallèles, aucune solution~; tout égal $\to$ même
droite, une infinité.

\textbf{9.} Les \omterm{def:g10:vectors:vector}{vecteurs} directeurs des droites sont $(-b, a)$ et
$(-d, c)$ (ou~: \omterm{def:g10:reffunc:affine}{coefficients directeurs} $-\frac ab$ et $-\frac cd$),
\omterm{def:g10:vectors:collinear}{colinéaires} exactement lorsque $ad - bc = 0$ --- le déterminant du
\cref{pb:g10:vectors:1} dans un nouvel emploi. Valeurs~: question 7,
$2 \times (-3) - 3 \times 5 = -21 \neq 0$~: solution unique.
Question 8~: $2 \times (-2) - (-1) \times 4 = 0$ dans les deux cas~:
parallèles ou confondues, et ce sont les constantes qui séparent les
deux cas.

\textbf{10.} $ad - bc = 6 - 2m$~: nul pour $m = 3$. Là,
$x + 3y = 3$ face à $2x + 6y = 7$~: en doublant la première on
obtient $2x + 6y = 6 \neq 7$, donc des droites parallèles et
\emph{aucune solution}. Pour $m \neq 3$, la combinaison linéaire
donne l'unique solution $y = \frac{1}{6 - 2m}$, puis
$x = 3 - \frac{m}{6 - 2m}$.

\textbf{11.} $p + r = 35$, $2p + 4r = 94$~: en divisant par $2$,
$p + 2r = 47$, donc $r = 12$ et $p = 23$~: vingt-trois faisans et
douze lapins --- la réponse même des \emph{Neuf Chapitres}.

\textbf{12.} $x + y = 30$ et $0.6x + 0.2y = 0.35 \times 30 = 10.5$~:
en retranchant $0.2 \times$ la première, $0.4x = 4.5$, donc
$x = 11.25$ L de sirop et $y = 18.75$ L de boisson.

\textbf{13.} Chiffres $a$ et $b$~: $a + b = 12$ et
$(10a + b) - (10b + a) = 9(a - b) = 18$, donc $a - b = 2$~:
$a = 7$, $b = 5$, et le nombre est $75$.

\textbf{14.} En additionnant~: $5c + 5r = 18$, donc $c + r = 3.60$
--- un café plus un croissant. En soustrayant~: $c - r = 0.20$.
D'où $c = 1.90$ et $r = 1.70$ euros~: le raccourci symétrique a
résolu le système sans la moindre substitution.

\textbf{15.} Doubler la première annonce donne $4$ pommes $+ 2$
poires $= 10$ euros, alors que l'étal en demande $9$~: le système est
incompatible --- droites parallèles, ensemble de solutions vide.
Verdict~: les deux annonces de prix ne peuvent pas être exactes
toutes les deux~; ou bien une remise se cache quelque part, ou bien
c'est l'arithmétique.

\textbf{16.} À l'origine~: $0 > 3 \times 0 - 1 = -1$, ce qui est
vrai~; donc $(0,0)$ appartient à la région $y > 3x - 1$ (au-dessus de
la droite). On teste n'importe quel point en le remplaçant dans
l'inéquation~: au-dessus ou au-dessous selon qu'elle est vraie ou
fausse.

\textbf{17.} Sommets~: $(0, 0)$~; $(4, 0)$ (axes et $x + y = 4$)~;
$(0, 1)$ (axe et $y = 2x + 1$)~; et $x + y = 4$ avec $y = 2x + 1$~:
$x = 1$, $y = 3$, soit $(1, 3)$.

\textbf{18.} $P(0,0) = 0$~; $P(4, 0) = 12$~; $P(0, 1) = 2$~;
$P(1, 3) = 9$. \omterm{def:g10:functions:extrema}{Maximum} au sommet $(4, 0)$~: le plan $x = 4$, $y = 0$
rapporte $12$ --- les sommets font la loi, comme le rendent visible
les droites parallèles $3x + 2y = c$ qui balaient la région.

\textbf{19.} $1.5x + 3 = 2x$ donne $x = 6$ et $y = 12$~: à six
kilomètres les deux taxis facturent douze euros --- le croisement du
devoir du week-end sur les deux thermomètres, du volume précédent,
devenu l'unique solution d'un système.

\textbf{20.} Les solutions d'une \omterm{def:g10:algebra:equation}{équation} du premier degré dessinent
une droite~; un système superpose deux droites, et $ad - bc$ dit d'un
coup d'œil si elles se coupent une fois (déterminant non nul) ou si
l'on tombe dans les cas parallèles-ou-confondues (déterminant nul).
Empilez plutôt des inéquations et les solutions forment une région
polygonale dont les sommets portent tout optimum linéaire. L'algèbre
linéaire généralise le déterminant à un nombre quelconque
d'inconnues~; la programmation linéaire industrialise les sommets.
\end{solution}
